\chapter{L'Idea: Unire Logica e Intelligenza Artificiale (L'Approccio Ibrido)}

Come abbiamo visto nel capitolo precedente, il vero problema dei moderni sistemi di 
riparazione automatica del codice (APR) non è la mancanza di ``creatività'' nel generare 
soluzioni, ma la mancanza di una \textbf{guida rigorosa}. Se lasciamo che i Large Language 
Model  facciano tutto da soli senza paletti rigidi, finiscono per inventarsi parti di codice inesistenti 
(allucinazioni) o creare soluzioni che funzionano solo in apparenza (overfitting). 
D'altra parte, gli strumenti puramente logici e formali sono dei chirurghi perfetti nel trovare 
\textit{dove} si trova l'errore, ma non hanno la flessibilità mentale per scrivere il codice che 
lo risolve in un contesto reale.

L'idea alla base di \textbf{HybridRepair} nasce proprio da qui: i due mondi, quello logico/simbolic
 e quello neurale/LLM, non sono nemici. Al contrario, sono \textbf{complementari}. Ognuno dei due è 
 fortissimo esattamente dove l'altro è debole.

\section{La Filosofia}

Il principio fondamentale di HybridRepair è una netta \textbf{separazione dei compiti}:

\begin{quote}
\textbf{L'analisi logica si occupa di trovare la causa del problema (diagnosi). L'LLM si occupa di scrivere il codice per risolverlo (sintesi).}
\end{quote}

Non si tratta di un semplice compromesso, ma di una scelta logica ben precisa dovuta alla natura diversa 
dei due compiti.

\textbf{La diagnosi:} Capire da dove arriva un valore \texttt{null}, seguire il 
percorso di un errore tra varie funzioni e beccare il punto esatto in cui il programma va in crash 
richiede un rigore logico e matematico assoluto. Sono garanzie che solo la programmazione logica 
(come Prolog) può offrire. Un LLM, per come è costruito, non può garantire questo livello 
di precisione deterministica.

\textbf{La sintesi:} Una volta che sappiamo con certezza scientifica \textit{cosa} 
correggere, il problema diventa \textit{come} scriverlo in un codice Java pulito, che rispetti lo 
stile del progetto e che faccia felici i test automatici. Qui gli LLM sono imbattibili: capiscono 
il contesto e scrivono codice flessibile in modo eccellente.

HybridRepair non chiede all'LLM di tirare a indovinare la causa del bug navigando alla cieca 
nel codice. E non chiede al motore logico di mettersi a scrivere codice Java. Ognuno fa il suo mestiere.

\section{L'Approccio Ibrido in Pratica: L'Esempio ``Closure-2''}

Per capire meglio, guardiamo un caso reale: il bug \textbf{Closure-2} (un errore di tipo 
\texttt{NullPointerException} o NPE) preso dal famoso database di test Defects4J, situato all'interno 
del Google Closure Compiler.

\subsection{Il problema}

Nel codice originale, c'è un'interfaccia che punta a un oggetto che non esiste. Quando il programma 
chiama il metodo \texttt{getImplicitPrototype()}, riceve in risposta  \texttt{null}. Alla riga 
successiva, il programma cerca di estrarre le proprietà da questo \texttt{null} usando \texttt{getOwnPropertyNames()}. 
Il risultato è che Il programma va in crash e la compilazione si blocca. I programmatori avevano persino lasciato 
un commento dicendo \textit{``questo è un errore di annotazione, ma non dovrebbe far crashare tutto''}, 
ma si erano dimenticati di inserire una protezione.

\subsection{Come fallisce un'IA ``pura'' (VibeRepair)}

Prendiamo VibeRepair \cite{zhu2026specification}, uno dei migliori sistemi APR su Defects4J basati solo su LLM del 2026. Prova a risolvere questo 
bug leggendo il codice e capendo (correttamente) che il metodo non deve andare in crash. Quindi, 
propone una via d'uscita rapida:

\begin{lstlisting}[language=Java]
if (implicitProto == null) {
    t.makeError(n, "Bad type annotation...");
    return; // <- L'IA decide di bloccare e uscire subito
}
\end{lstlisting}

Sembra un'ottima soluzione: il crash sparisce e i test passano. Ma \textbf{è un errore concettuale}. 
Facendo \texttt{return}, l'IA interrompe bruscamente il lavoro. Poche righe più in basso c'è un ciclo 
\texttt{for-each} che dovrebbe analizzare altre interfacce valide, ma a causa di questa ``uscita anticipata'' 
verranno tutte ignorate in silenzio.

VibeRepair capisce che il programma deve sopravvivere, ma non si accorge del ciclo sottostante che ci 
dice chiaramente quale sia la vera soluzione: \textbf{passare un insieme vuoto}, in modo che il ciclo faccia 
``zero giri'' senza interrompere il resto del programma.

\subsection{Come vince HybridRepair}

Ecco come l' approccio ibrido risolve la situazione in tre passi:

\begin{enumerate}
    \item \textbf{Step 1 (LogicFL).} Il sistema logico fa girare i test falliti e 
    ricostruisce in modo inequivocabile la catena del disastro:
\begin{lstlisting}[language=Prolog]
null_origin(implicitProto, line(1570), method('getImplicitPrototype')).
npe_cause(implicitProto, line(1571), method('getOwnPropertyNames')).
\end{lstlisting}
    L'LLM non si sente dire ``Ehi, c'è un crash da qualche parte''. Riceve un rapporto precisissimo: 
    \textit{``La variabile \texttt{implicitProto} nasce \texttt{null} alla riga 1570 e fa esplodere tutto alla riga 1571''}.
    
    \item \textbf{Step 2 (Estrazione AST).} \\
    Il sistema isola chirurgicamente solo le $\sim$30 righe del metodo difettoso, scartando le altre 1700 
    righe del file. Così facendo, l'LLM ha una visuale chiara e vede immediatamente il ciclo \texttt{for-each} incriminato.
    
    \item \textbf{Step 3  (Sintesi LLM).} \\
    Ora l'LLM ha tutto: la causa esatta e il contesto pulito. Ragiona sul problema e capisce che l'obiettivo 
    è non far crashare la riga 1571, ma permettere al ciclo successivo di continuare. La soluzione è 
    inizializzare la variabile con una lista vuota (\texttt{HashSet}):
\begin{lstlisting}[language=Java]
// Rimossa la riga originale:
% currentPropertyNames = implicitProto.getOwnPropertyNames();
// Inserita la patch correttiva:
if (implicitProto == null) {
    currentPropertyNames = new HashSet<>();
} else {
    currentPropertyNames = implicitProto.getOwnPropertyNames();
}
\end{lstlisting}
\end{enumerate}

Questa patch fa esattamente la stessa cosa che hanno fatto gli ingegneri di Google nella realtà 
(loro hanno usato \texttt{ImmutableSet.of()}, ma la logica è identica). I test passano e, soprattutto, 
il comportamento del software rimane corretto.

\subsection{Il Confronto}

\begin{table}[h]
\centering
\caption{Confronto tra VibeRepair e HybridRepair}
\vspace{2mm}
\begin{tabular}{lll}
\toprule
\textbf{Caratteristica} & \textbf{VibeRepair (Solo IA)} & \textbf{HybridRepair (Logica + IA)} \\
\midrule
Ricerca del guasto & Ha bisogno di aiuti esterni & Completamente autonomo (LogicFL) \\
Causa dell'errore & Tirata a indovinare dal codice & Dimostrata matematicamente (Prolog) \\
Cosa vede l'IA & Un'idea vaga del problema & Rapporto esatto + solo il codice utile \\
Codice generato & \texttt{return} anticipato (Sbagliato) & Protezione con \texttt{HashSet} vuoto (Corretto) \\
Supera i Test? & Sì  & Sì  \\
È davvero corretto? & No ($\times$) & Sì  \\
\bottomrule
\end{tabular}
\end{table}

Questo esempio dimostra una cosa fondamentale: la differenza tra una patch che ``sembra giusta'' e 
una che ``è giusta'' non dipende da quanto è intelligente l'IA, ma dalla \textbf{qualità delle informazioni} 
che le diamo in pasto.

\section{Come è costruito: I Tre Livelli}

Tutto questo discorso si traduce in un'architettura software divisa in tre fasi, come una vera e
 propria catena di montaggio:

\begin{figure}[h]
\centering
\begin{tikzpicture}[
    node distance=1.2cm,
    block/.style={rectangle, draw, fill=gray!10, text width=4cm, align=center, rounded corners, minimum height=1.5cm, font=\small},
    arrow/.style={thick,->,>=stealth}
]
\node (oracle) [block] {\textbf{1. FaultOracle}\\(L'Oracolo dei Guasti)\\Analisi puramente logica};
\node (reason) [block, right=0.8cm of oracle] {\textbf{2. SpecReason}\\(Il Ragionatore)\\Ponte Logica-Codice};
\node (agent) [block, right=0.8cm of reason] {\textbf{3. PatchAgent}\\(L'Agente Programmatore)\\Sintesi e Validazione};

\draw [arrow] (oracle) -- (reason);
\draw [arrow] (reason) -- (agent);
\end{tikzpicture}
\caption{Pipeline a tre livelli di HybridRepair.}
\end{figure}

\begin{enumerate}
    \item \textbf{FaultOracle (L'Oracolo dei Guasti):} È la parte puramente logica (Prolog). Niente IA, 
            niente probabilità. Trova la causa del bug in modo matematico e fornisce un referto inequivocabile.
    \item \textbf{SpecReason (Il Ragionatore):} È il ponte tra la logica e il codice. Qui l'LLM entra in gioco 
            ma \textit{non scrive ancora codice}. Viene obbligato a pensare ad alta voce (Chain-of-Thought) 
            rispondendo a tre domande: \textit{Cosa c'è di sbagliato? Cosa vogliamo ottenere? Come lo risolviamo 
            a livello di concetto?} Genera comprensione, non codice.
    \item \textbf{PatchAgent (L'Agente Programmatore):} È la fase finale. L'LLM, partendo dai ragionamenti 
            dello step precedente e avendo il contesto pulito, scrive finalmente la patch. La applica, prova 
            a compilarla e, se qualcosa va storto, impara dall'errore e ci riprova senza fare due volte lo stesso sbaglio.
\end{enumerate}

Questa struttura a tre strati è il segreto di HybridRepair: l'IA non viene mai usata per investigare sulle 
cause (dove farebbe danni), ma viene usata solo per tradurre in codice delle verità già dimostrate dalla logica.